% !TEX program = xelatex
% ============================================================
%  Deep Generative Models — Chapter 3: Denoising Diffusion Probabilistic Models
% ============================================================

\documentclass[11pt,a4paper]{article}

% ============================================================
% §0  Fonts
% ============================================================
% ctex is only needed for CJK typesetting; load it only when available.
\IfFileExists{ctex.sty}{\usepackage{ctex}}{}

% ============================================================
% §1  Page layout
% ============================================================
\IfFileExists{geometry.sty}{%
  \usepackage[margin=2.5cm]{geometry}
}{%
  \PackageWarning{geometry}{geometry package not found; using default A4 margins. Install `geometry' for 2.5cm margins.}%
}
\usepackage{microtype}
\linespread{1.08}

% ============================================================
% §2  Packages
% ============================================================
\usepackage{amsmath,amssymb,amsthm}
\renewcommand{\proofname}{Proof}
\usepackage{mathtools}
\usepackage{graphicx}
\graphicspath{{../graph/}}
\usepackage{float}
\usepackage{enumitem}
\usepackage{soul}
\usepackage{leftindex}
\usepackage{tikz}
\usetikzlibrary{arrows,calc,positioning}
\usepackage[dvipsnames]{xcolor}
\usepackage[most]{tcolorbox}
\usepackage{fontawesome5}
\usepackage{titlesec}

% ============================================================
% §3  Colors
% ============================================================
\definecolor{mydarkblue}{rgb}{0.0,0.08,0.45}
\definecolor{nicepink}  {RGB}{255,105,180}
\definecolor{cardbg}    {RGB}{245,245,247}
\definecolor{blue_new}  {RGB}{63,77,238}
\definecolor{page_white}{RGB}{227,227,229}
\definecolor{highlight} {HTML}{FCD66F}
\definecolor{warning}   {HTML}{E57373}
\definecolor{info}      {HTML}{4DB6AC}
\definecolor{navblue}     {HTML}{0077b6}
\definecolor{navcyan}     {HTML}{00b4d8}
\definecolor{navdark}     {HTML}{023e8a}
\definecolor{accentorange}{HTML}{fb8500}

\newcommand{\red}   [1]{{\color{OrangeRed}#1}}
\newcommand{\tlight}[1]{\textcolor{warning}{#1}}
\newcommand{\tinfo} [1]{\textcolor{info}{#1}}
\newcommand{\bglight}[1]{\hl{#1}}
\newcommand{\cbox}[2]{\colorbox{#1}{\textcolor{black}{#2}}}
\newcommand{\bgwarn}[1]{\cbox{warning}{#1}}
\newcommand{\bginfo}[1]{\cbox{info}{#1}}

% ============================================================
% §4  Hyperlinks
% ============================================================
\usepackage{hyperref}
\hypersetup{
    colorlinks = true,
    linkcolor  = mydarkblue,
    citecolor  = mydarkblue,
    urlcolor   = mydarkblue
}

% ============================================================
% §5  Math macros
% ============================================================
\newcommand{\wh}{\widehat}
\newcommand{\wt}{\widetilde}
\newcommand{\ov}{\overline}
\renewcommand{\tilde}{\wt}
\renewcommand{\hat}{\wh}
\renewcommand{\d}{\mathrm{d}}

\newcommand{\N}{\mathbb{N}}  \newcommand{\R}{\mathbb{R}}
\newcommand{\Q}{\mathbb{Q}}  \newcommand{\Z}{\mathbb{Z}}
\newcommand{\C}{\mathbb{C}}

\DeclareMathOperator*{\E}{\mathbb{E}}
\let\P\relax
\DeclareMathOperator*{\P}{\mathbb{P}}
\DeclareMathOperator*{\var}{Var}
\DeclareMathOperator*{\argmin}{arg\,min}
\DeclareMathOperator*{\argmax}{arg\,max}
\DeclareMathOperator{\diag}{diag}
\DeclareMathOperator{\tr}{tr}
\DeclareMathOperator{\rank}{rank}
\DeclareMathOperator{\dist}{dist}

\renewcommand{\L}{\mathcal{L}}
\newcommand{\cN}{\mathcal{N}}

% ---- DDPM-specific ----
\newcommand{\KL}[2]{D_{\mathrm{KL}}\!\left(#1\,\|\,#2\right)}
\newcommand{\balpha}{\bar{\alpha}}
\newcommand{\noiseschedule}{\{\beta_i\}_{i=1}^{L}}
\newcommand{\pdata}{p_{\mathrm{data}}}
\newcommand{\pprior}{p_{\mathrm{prior}}}
\newcommand{\Ldiff}{\mathcal{L}_{\text{diffusion}}}
\newcommand{\Lprior}{\mathcal{L}_{\text{prior}}}
\newcommand{\Lrecon}{\mathcal{L}_{\text{recon}}}
\newcommand{\Lsimple}{\mathcal{L}_{\text{simple}}}
\newcommand{\LDDPM}{\mathcal{L}_{\text{DDPM}}}
\newcommand{\LELBO}{\mathcal{L}_{\text{ELBO}}}
\newcommand{\Lbound}{\mathcal{B}}
% \Efwd: shorthand for expectation over forward trajectory x_{1:L} given fixed x_0
% NOTE: this macro is ONLY used inside the ELBO proof after x_0 is fixed;
%       do NOT redefine it inline (that would produce a tautology).
\newcommand{\Efwd}{\mathbb{E}^{\mathrm{fwd}}_{x_0}}

% p_i(x_i|x_0): forward law given x_0; \bar p_i(x_i): population marginal over x_0\sim p_data
\newcommand{\bpMarg}[1]{\bar p_{#1}}

% ============================================================
% §6  Section heading style
% ============================================================
\titleformat{\section}
  {\large\bfseries\color{mydarkblue}}{\thesection}{0.5em}{}[\titlerule]
\titlespacing*{\section}{0pt}{1.5ex}{1ex}

\titleformat{\subsection}
  {\normalsize\bfseries}{\thesubsection}{0.5em}{}
\titlespacing*{\subsection}{0pt}{1ex}{0.5ex}

% ============================================================
% §7  Compact title + horizontal TOC
% ============================================================
\makeatletter
\newcommand{\makecompacttitle}{%
    \begin{tcolorbox}[
        colback=mydarkblue!5, colframe=mydarkblue!10,
        arc=3pt, boxrule=0.5pt,
        left=5pt, right=5pt, top=5pt, bottom=5pt
    ]
        {\large\bfseries \@title}\hfill{\small \@author}\\[-2pt]
        {\tiny\color{gray} \@date \hfill Deep Generative Models Notes}
    \end{tcolorbox}
    \vspace{-1em}
}
\makeatother

\newtcolorbox{tabNavBoxSmall}[1][]{
    enhanced,
    colframe=white, colback=gray!5, coltext=black!80,
    fontupper=\sffamily\small\bfseries,
    arc=1.5mm, boxrule=0pt, bottomrule=2pt,
    left=2mm, right=2mm, top=1.5mm, bottom=1.5mm,
    #1
}

\newcommand{\horizontalTOCDDPM}{%
    \begin{center}
    \begin{tcbraster}[
        raster columns=2, raster width=\linewidth,
        raster column skip=1mm, raster row skip=1mm, raster equal height
    ]
        \begin{tabNavBoxSmall}\hyperref[sec:overview]{\ Overview}\end{tabNavBoxSmall}
        \begin{tabNavBoxSmall}\hyperref[sec:forward]{\ Forward Process}\end{tabNavBoxSmall}
        \begin{tabNavBoxSmall}\hyperref[sec:reverse]{\ Reverse Process}\end{tabNavBoxSmall}
        \begin{tabNavBoxSmall}\hyperref[sec:elbo]{\ DDPM ELBO}\end{tabNavBoxSmall}
        \begin{tabNavBoxSmall}\hyperref[sec:practical]{\ Practical Choices}\end{tabNavBoxSmall}
        \begin{tabNavBoxSmall}\hyperref[sec:sampling]{\ Sampling}\end{tabNavBoxSmall}
    \end{tcbraster}
    \vspace{-2.5pt}%
    \noindent\textcolor{gray!30}{\rule{\linewidth}{1pt}}
    \end{center}
    \vspace{1em}
}

% ============================================================
% §8  Theorem environments (tcolorbox)
% ============================================================
\tcbset{
  mytheo/.style={
    enhanced, breakable,
    colback=#1!5!white, colframe=#1!70!black,
    coltitle=white, fonttitle=\bfseries\small,
    boxrule=0.6pt, arc=3pt,
    left=5pt, right=5pt, top=10pt, bottom=5pt,
    attach boxed title to top left={yshift*=-\tcboxedtitleheight/2, xshift=2mm},
    boxed title style={
      colback=#1!80!black, colframe=#1!80!black,
      boxrule=0pt, arc=2pt
    }
  }
}

\newtcbtheorem[auto counter]{theorem}    {Theorem}{mytheo=WildStrawberry}{thm}
\newtcbtheorem[auto counter]{definition} {Definition}{mytheo=NavyBlue  }{def}
\newtcbtheorem[auto counter]{example}    {Example}{mytheo=YellowOrange }{ex}
\newtcbtheorem[auto counter]{remark}     {Remark}{mytheo=Periwinkle    }{rem}
\newtcbtheorem[auto counter]{proposition}{Proposition}{mytheo=JungleGreen}{prop}
\newtcbtheorem[auto counter]{lemma}      {Lemma}{mytheo=ForestGreen  }{lem}
\newtcbtheorem[auto counter]{corollary}  {Corollary}{mytheo=BrickRed   }{cor}
\newtcbtheorem[auto counter]{fact}       {Fact}{mytheo=YellowGreen  }{fact}
\newtcbtheorem[auto counter]{assumption} {Assumption}{mytheo=Plum     }{as}

\tcolorboxenvironment{proof}{
  blanker, breakable, left=4pt, right=4pt, top=3pt, bottom=3pt,
  borderline west={2pt}{0pt}{gray!50}
}

% ============================================================
% §9  Body
% ============================================================
\title{Chapter 3: Denoising Diffusion Probabilistic Models}
\author{}
\date{\today}

\begin{document}

\makecompacttitle
\horizontalTOCDDPM

\begin{center}
  \includegraphics[width=\linewidth]{\detokenize{Gemini_Generated_Image_nxh0finxh0finxh0.png}}
\end{center}
\vspace{0.5em}

% ============================================================
\section{Overview}\label{sec:overview}
% ============================================================

\begin{definition}{DDPM as Two Stochastic Processes}{ddpm-two-processes}
A denoising diffusion probabilistic model (DDPM) is defined by two coupled stochastic processes:
\begin{itemize}[leftmargin=1.2em, itemsep=2pt]
    \item a \textbf{forward process} that gradually corrupts data by Gaussian noise;
    \item a \textbf{reverse denoising process} that learns to remove noise step by step.
\end{itemize}
\end{definition}

The forward chain starts at real data $x_0\sim\pdata$ and moves toward an isotropic Gaussian prior:
\[
x_0 \rightarrow x_1 \rightarrow \cdots \rightarrow x_{L},\qquad
x_{L}\sim \bpMarg{L},\quad \bpMarg{L}\approx \pprior=\cN(0,\mathbf{I}).
\]
The reverse chain learns a parameterized transition
\[
q_{\phi}(x_{i-1}\mid x_i),\qquad i=L,\dots,1,
\]
to transform unstructured Gaussian noise back to structured samples.

% ============================================================
\section{Forward Process (Fixed Encoder)}\label{sec:forward}
% ============================================================

\subsection{Markov Corruption Kernel}

\begin{definition}{Forward Transition}{forward-def}
Let $\noiseschedule$ be a predefined monotone increasing schedule. Define
\[
\alpha_i \triangleq \sqrt{1-\beta_i},\qquad i=1,\dots,L.
\]
The forward process is the Markov chain
\[
p(x_i\mid x_{i-1})
=\cN\!\left(x_i;\,\alpha_i x_{i-1},\,\beta_i\mathbf{I}\right).
\]
\end{definition}

\begin{remark}{Noise Schedule Intuition}{forward-remark}
Each forward step first scales the previous state by $\alpha_i$ and then injects isotropic Gaussian noise with variance $\beta_i$. As $i$ grows, signal shrinks and noise accumulates.
\end{remark}

\subsection{Closed-Form Marginal at Arbitrary Time}

Define the cumulative product
\[
\balpha_i \triangleq \prod_{k=1}^{i}\alpha_k
= \prod_{k=1}^{i}\sqrt{1-\beta_k}.
\]

\begin{proposition}{Closed Form of \texorpdfstring{$p_i(x_i\mid x_0)$}{pi(xi|x0)}}{forward-closed}
For every $i\in\{1,\dots,L\}$, the forward law at noise level~$i$ given clean $x_0$ is
\[
p_i(x_i\mid x_0)
=\cN\!\left(x_i;\,\balpha_i x_0,\,(1-\balpha_i^2)\mathbf{I}\right).
\]
Equivalently,
\[
x_i = \balpha_i x_0 + \sqrt{1-\balpha_i^2}\,\varepsilon,\qquad
\varepsilon\sim\cN(0,\mathbf{I}).
\]
\end{proposition}

\begin{proof}
Use induction on $i$. The statement is immediate for $i=1$. Assume it holds for $i-1$. Since
\[
x_i=\alpha_i x_{i-1}+\sqrt{\beta_i}\,\xi_i,\qquad \xi_i\sim\cN(0,\mathbf{I}),
\]
substitute the induction form of $x_{i-1}$:
\[
x_i=\alpha_i\!\left(\balpha_{i-1}x_0+\sqrt{1-\balpha_{i-1}^2}\,\varepsilon'\right)+\sqrt{\beta_i}\xi_i
=\balpha_i x_0+\eta_i,
\]
where $\eta_i$ is Gaussian with covariance
\[
\alpha_i^2(1-\balpha_{i-1}^2)\mathbf{I}+\beta_i\mathbf{I}
=(1-\balpha_i^2)\mathbf{I}.
\]
Thus $x_i\mid x_0$ is Gaussian with the claimed mean and covariance.
\end{proof}

The forward marginal $p_i(x_i\mid x_0)=\cN(\bar\alpha_i x_0,(1-\bar\alpha_i^2)\mathbf{I})$
converges to $\cN(0,\mathbf{I})$ and becomes independent of $x_0$ if and only if
$\bar\alpha_i=\prod_{k=1}^{i}\sqrt{1-\beta_k}\to 0$.
A monotone-increasing schedule $\{\beta_i\}$ bounded away from zero is a sufficient (but not necessary) condition for $\bar\alpha_L\to 0$.
Integrating over the unknown data law $\pdata$ gives the population marginal
\[
\bpMarg{i}(x_i)=\int \pdata(x_0)\,p_i(x_i\mid x_0)\,\d x_0,
\]
\begin{remark}{Marginal Levels and Kernel Naming}{marginal-kernel-naming}
The population marginal differs at each noise level
($\bpMarg{i}\neq \bpMarg{i'}$ for $i\neq i'$), even though every one-step kernel
$p(x_i\mid x_{i-1})$ is known in closed form.
We write $\pdata$ for the clean-data distribution; $p_i(x_i\mid x_0)$ for the
forward law of $x_i$ given~$x_0$; and
$\bpMarg{i}(x_i)=\int \pdata(x_0)\,p_i(x_i\mid x_0)\,\d x_0$ for its population
marginal.
Only adjacent-time maps $p(x_i\mid x_{i-1})$ are called \emph{transition
kernels}. Reverse \emph{kernels} are denoted by~$q$ (learned transitions use
$q_\phi$).
This is the opposite of the original DDPM paper \emph{(Ho et al.\ 2020)}, which
uses $q$ for the forward process and $p_\theta$ for the reverse model.
\end{remark}

% ============================================================
\section{Reverse Denoising Process}\label{sec:reverse}
% ============================================================

\subsection{Learning an Approximate Reverse Kernel}

The true reverse transition $q(x_{i-1}\mid x_i)$ is unknown.
DDPM learns a tractable Gaussian family:
\[
q_\phi(x_{i-1}\mid x_i)
\triangleq \cN\!\left(x_{i-1};\,\mu_\phi(x_i,i),\,\sigma^2(i)\mathbf{I}\right),
\]
where $\mu_\phi(\cdot,i):\R^D\to\R^D$ is learnable and $\sigma^2(i)>0$ is a fixed
schedule.  Following the original DDPM paper, $\sigma^2(i)$ is set equal to the
forward-posterior variance
$\sigma^2(i)=\frac{(1-\bar\alpha_{i-1}^2)\beta_i}{1-\bar\alpha_i^2}$
(derived in Proposition~\ref{prop:reverse-posterior}), which ensures the
mean-matching KL identity used throughout.

The stage-wise objective is to minimize
\[
\E_{x_i\sim \bpMarg{i}}\!\left[\KL{q(x_{i-1}\mid x_i)}{q_\phi(x_{i-1}\mid x_i)}\right].
\]

\subsection{Conditioning on \texorpdfstring{$x_0$}{x0}}

By Bayes' rule under the forward chain,
\[
q(x_{i-1}\mid x_i)
=\frac{p(x_i\mid x_{i-1})\,\bpMarg{i-1}(x_{i-1})}{\bpMarg{i}(x_i)}.
\]
Each forward kernel $p(x_i\mid x_{i-1})$ is explicit, but the reverse kernel still depends on the marginals $\bpMarg{i-1}$ (or $\pdata$ when $i=1$) and $\bpMarg{i}$.
The latter is obtained from $\pdata$ via
\[
\bpMarg{i}(x_i)=\int \pdata(x_0)\,p_i(x_i\mid x_0)\,\d x_0,
\]
and $\pdata$ is unknown in practice.
Therefore, even with a fully specified forward process, we cannot evaluate $\bpMarg{i}(x_i)$ or $q(x_{i-1}\mid x_i)$ pointwise.
The remedy is to condition on a clean sample $x_0\sim\pdata$, for which all forward quantities are Gaussian:
\[
p(x_{i-1}\mid x_i,x_0)
=p(x_i\mid x_{i-1})\frac{p_{i-1}(x_{i-1}\mid x_0)}{p_i(x_i\mid x_0)}.
\]
All three factors on the right are Gaussian, so the posterior above is also Gaussian.
Training then replaces expectations under $\bpMarg{i}$ by Monte Carlo over $x_0\sim\pdata$ and $x_i\sim p_i(\cdot\mid x_0)$.

\begin{theorem}[breakable=false]{KL Decomposition for Reverse Matching}{reverse-kl-decomp}
For each $i$, define the marginal and posterior reverse-KL objectives
\begin{align*}
\mathcal{J}^{\mathrm{rev}}_i(\phi)
&\triangleq \E_{x_i\sim \bpMarg{i}}\!\left[\KL{q(x_{i-1}\mid x_i)}{q_\phi(x_{i-1}\mid x_i)}\right],\\
\mathcal{J}^{\mathrm{post}}_i(\phi)
&\triangleq \E_{x_0\sim\pdata}\,\E_{x_i\sim p_i(\cdot\mid x_0)}
\!\left[\KL{p(x_{i-1}\mid x_i,x_0)}{q_\phi(x_{i-1}\mid x_i)}\right].
\end{align*}
Then
\[
\mathcal{J}^{\mathrm{rev}}_i(\phi)=\mathcal{J}^{\mathrm{post}}_i(\phi)+C_i,
\]
where $C_i$ is independent of $\phi$.
In particular, minimizing the conditional objective is equivalent to minimizing the marginal reverse-KL up to an additive constant.
\end{theorem}

\begin{proof}
Fix $x_i$. The true reverse kernel is a mixture of tractable posteriors over $x_0$:
\[
q(x_{i-1}\mid x_i)
=\int p(x_{i-1}\mid x_i,x_0)\,\pdata(x_0\mid x_i)\,\d x_0
=\E_{x_0\sim \pdata(\cdot\mid x_i)}\!\left[p(x_{i-1}\mid x_i,x_0)\right],
\]
where $\pdata(x_0\mid x_i)\propto \pdata(x_0)\,p_i(x_i\mid x_0)$.
Subtract the conditional objective from the marginal KL; the $q_\phi$ terms cancel:
\begin{align*}
&\E_{x_i\sim \bpMarg{i}}\!\left[\KL{q(x_{i-1}\mid x_i)}{q_\phi(x_{i-1}\mid x_i)}\right]
-\E_{x_0\sim\pdata}\,\E_{x_i\sim p_i(\cdot\mid x_0)}
\!\left[\KL{p(x_{i-1}\mid x_i,x_0)}{q_\phi(x_{i-1}\mid x_i)}\right]\\
&\quad=
\E_{(x_{i-1},x_i)\sim p(x_{i-1},x_i)}\!\left[\log q(x_{i-1}\mid x_i)\right]
-\E_{(x_0,x_{i-1},x_i)\sim p(x_0,x_{i-1},x_i)}\!\left[\log p(x_{i-1}\mid x_i,x_0)\right]\\
&\quad=
\E_{x_i\sim \bpMarg{i}}
\E_{x_0\sim \pdata(\cdot\mid x_i)}
\Big[
\log \pdata(x_0\mid x_i)
-\E_{x_{i-1}\sim p(\cdot\mid x_i,x_0)}\!\left[\log \pdata(x_0\mid x_{i-1},x_i)\right]
\Big],
\end{align*}
using $\log q(x_{i-1}\mid x_i)-\log p(x_{i-1}\mid x_i,x_0)
=\log \pdata(x_0\mid x_i)-\log \pdata(x_0\mid x_{i-1},x_i)$,
where $p(x_{i-1},x_i)=p(x_i\mid x_{i-1})\,\bpMarg{i-1}(x_{i-1})$ and
$p(x_0,x_{i-1},x_i)=\pdata(x_0)\,p(x_{i-1}\mid x_i,x_0)\,p_i(x_i\mid x_0)$.
This depends only on the fixed forward chain, hence it is a $\phi$-free constant $C_i$.
Therefore $\mathcal{J}^{\mathrm{rev}}_i(\phi)=\mathcal{J}^{\mathrm{post}}_i(\phi)+C_i$.
\end{proof}

\begin{remark}{Why Conditional Training Is Legitimate (Two Independent Routes)}{conditional-training-legitimacy}
The stage-wise goal is to match the \emph{marginal} reverse kernel $q(x_{i-1}\mid x_i)$, but this object depends on unknown marginals $\bpMarg{i-1},\bpMarg{i}$ and cannot be evaluated directly.
Conditioning on a clean sample $x_0$ replaces it with the tractable posterior $p(x_{i-1}\mid x_i,x_0)$.

\textbf{Route~A (Theorem~1 only).}
Theorem~\ref{thm:reverse-kl-decomp} shows
$\mathcal{J}^{\mathrm{rev}}_i(\phi)=\mathcal{J}^{\mathrm{post}}_i(\phi)+C_i$ with $C_i$ independent of~$\phi$.
So minimizing the conditional objective is exactly the same as minimizing the marginal reverse-KL up to an additive constant.
This route does \emph{not} require the ELBO of Section~\ref{sec:elbo}.

\textbf{Route~B (original DDPM paper).}
The paper instead starts from maximizing $\log q_\phi(x_0)$ and derives a variational upper bound (Theorem~\ref{thm:ddpm-elbo}), whose diffusion term already has the conditional form $\KL{p(x_{i-1}\mid x_i,x_0)}{q_\phi(x_{i-1}\mid x_i)}$.
That proof is self-contained and does not use Theorem~\ref{thm:reverse-kl-decomp}.

Both routes justify the same practical recipe---sample $x_0\sim\pdata$, corrupt to $x_i$, and train $q_\phi$ against the Gaussian forward posterior---but they answer the question at different levels: Route~A is step-wise reverse matching; Route~B is end-to-end likelihood maximization.
\end{remark}

\begin{proposition}{Closed Form of \texorpdfstring{$p(x_{i-1}\mid x_i,x_0)$}{p(xi-1|xi,x0)}}{reverse-posterior}
The conditional posterior is
\[
p(x_{i-1}\mid x_i,x_0)
=\cN\!\left(x_{i-1};\,\mu(x_i,x_0,i),\,\sigma^2(i)\mathbf{I}\right),
\]
with
\[
\mu(x_i,x_0,i)
=\frac{\alpha_i-\balpha_i^2/\alpha_i}{1-\balpha_i^2}x_i
+\frac{(1-\alpha_i^2)\balpha_i/\alpha_i}{1-\balpha_i^2}x_0,
\]
\[
\sigma^2(i)
=\frac{1-\balpha_{i-1}^2}{1-\balpha_i^2}\,\beta_i.
\]
\end{proposition}

\begin{proof}
This is the adjacent-step special case of the general reverse conditional
transition formula for linear Gaussian corruptions; the general case for
arbitrary $0\le t<s$ will be discussed later.
We prove the one-step DDPM case directly.

Conditioned on $x_0$, the forward marginal and one-step transition are
\[
p_{i-1}(x_{i-1}\mid x_0)
=\cN\!\left(x_{i-1};\,\balpha_{i-1}x_0,\,(1-\balpha_{i-1}^2)\mathbf{I}\right),
\qquad
p(x_i\mid x_{i-1})
=\cN\!\left(x_i;\,\alpha_i x_{i-1},\,\beta_i\mathbf{I}\right).
\]
By Bayes' rule,
\[
p(x_{i-1}\mid x_i,x_0)
\propto
p(x_i\mid x_{i-1})\,p_{i-1}(x_{i-1}\mid x_0).
\]
As a function of $x_{i-1}$, the log-density is quadratic, so the posterior is
Gaussian.
Its precision is
\[
\frac{\alpha_i^2}{\beta_i}
+\frac{1}{1-\balpha_{i-1}^2}
=\frac{1-\balpha_i^2}{\beta_i(1-\balpha_{i-1}^2)},
\]
hence
\[
\sigma^2(i)
=\frac{\beta_i(1-\balpha_{i-1}^2)}{1-\balpha_i^2}.
\]
The corresponding mean is
\[
\mu(x_i,x_0,i)
=\sigma^2(i)\!\left(
\frac{\alpha_i}{\beta_i}x_i
+\frac{\balpha_{i-1}}{1-\balpha_{i-1}^2}x_0
\right)
=\frac{\alpha_i(1-\balpha_{i-1}^2)}{1-\balpha_i^2}x_i
+\frac{\beta_i\balpha_{i-1}}{1-\balpha_i^2}x_0.
\]
Since $\balpha_{i-1}=\balpha_i/\alpha_i$ and $\beta_i=1-\alpha_i^2$, this is
equivalently
\[
\mu(x_i,x_0,i)
=\frac{\alpha_i-\balpha_i^2/\alpha_i}{1-\balpha_i^2}x_i
+\frac{(1-\alpha_i^2)\balpha_i/\alpha_i}{1-\balpha_i^2}x_0,
\]
which proves the claimed form.
\end{proof}

This yields the diffusion training term (mean-matching form,
already averaged over $x_0$):
\[
\Ldiff(\phi)
\triangleq
\sum_{i=2}^{L}
\frac{1}{2\sigma^2(i)}
\E_{x_0\sim\pdata}\,\E_{x_i\sim p_i(\cdot\mid x_0)}
\!\left[
\left\|
\mu_\phi(x_i,i)-\mu(x_i,x_0,i)
\right\|_2^2
\right].
\]
Because $q_\phi$ and $p(x_{i-1}\mid x_i,x_0)$ share the same variance $\sigma^2(i)$,
the Gaussian KL satisfies
$\KL{p(x_{i-1}\mid x_i,x_0)}{q_\phi(x_{i-1}\mid x_i)}
=\frac{1}{2\sigma^2(i)}\|\mu(x_i,x_0,i)-\mu_\phi(x_i,i)\|_2^2$,
so $\Ldiff(\phi)=\E_{x_0\sim\pdata}[\Ldiff(x_0;\phi)]$ where $\Ldiff(x_0;\phi)$ is the
per-sample KL sum appearing in the ELBO (Section~\ref{sec:elbo}).

% ============================================================
\section{DDPM ELBO}\label{sec:elbo}
% ============================================================

\subsection{Joint Model and Variational Bound}

Define the latent trajectory $x_{1:L}$ and generative joint
\[
q_\phi(x_0,x_{1:L})
=\pprior(x_L)\prod_{i=1}^{L}q_\phi(x_{i-1}\mid x_i),
\]
where $\pprior(x_L)=\cN(0,\mathbf{I})$ is the sampling prior at the top of the reverse chain (the forward marginal satisfies $\bpMarg{L}\approx\pprior$ by schedule design).
Let the forward trajectory given $x_0$ be Markov with transitions $p(x_i\mid x_{i-1})$, so
\[
\prod_{i=1}^{L}p(x_i\mid x_{i-1})
\]
is the joint density of $(x_1,\dots,x_L)$ and $\bpMarg{i}(x_i)=\int \pdata(x_0)\,p_i(x_i\mid x_0)\,\d x_0$.

\begin{theorem}{ELBO Decomposition for DDPM}{ddpm-elbo}
For any data sample $x_0$,
\[
-\log q_\phi(x_0)
\le -\LELBO(x_0;\phi)
=\Lprior(x_0)+\Lrecon(x_0;\phi)+\Ldiff(x_0;\phi),
\]
with
\[
\Lprior(x_0)
\triangleq \KL{p_L(\cdot\mid x_0)}{\pprior(x_L)},
\]
\[
\Lrecon(x_0;\phi)
\triangleq \E_{x_1\sim p_1(\cdot\mid x_0)}\!\left[-\log q_\phi(x_0\mid x_1)\right],
\]
\[
\Ldiff(x_0;\phi)
\triangleq \sum_{i=2}^{L}
\E_{x_i\sim p_i(\cdot\mid x_0)}
\!\left[
\KL{p(x_{i-1}\mid x_i,x_0)}{q_\phi(x_{i-1}\mid x_i)}
\right].
\]
\end{theorem}

\begin{proof}
Fix $x_0$. The forward joint factorises as
\[
p(x_{1:L}\mid x_0)=\prod_{i=1}^{L}p(x_i\mid x_{i-1}).
\]
Write $\Efwd[\cdot]\triangleq\E_{x_{1:L}\sim p(x_{1:L}\mid x_0)}[\cdot]$
for the expectation over the full forward trajectory given this fixed~$x_0$.
Then
\begin{align*}
-\log q_\phi(x_0)
&= -\log \int q_\phi(x_0,x_{1:L})\,\d x_{1:L} \\
&= -\log \int p(x_{1:L}\mid x_0)\,
\frac{q_\phi(x_0,x_{1:L})}{p(x_{1:L}\mid x_0)}\,\d x_{1:L} \\
&= -\log \Efwd\!\left[\frac{q_\phi(x_0,x_{1:L})}{p(x_{1:L}\mid x_0)}\right] \\
&\le \Efwd\!\left[-\log q_\phi(x_0,x_{1:L})
+\sum_{i=1}^{L}\log p(x_i\mid x_{i-1})\right]
\triangleq \Lbound(x_0;\phi),
\end{align*}
by Jensen's inequality ($-\log$ is convex).
Using
$q_\phi(x_0,x_{1:L})=\pprior(x_L)\prod_{i=1}^{L}q_\phi(x_{i-1}\mid x_i)$,
expand the bound:
\begin{align*}
\Lbound(x_0;\phi)
&= \Efwd\!\left[-\log \pprior(x_L)+\sum_{i=1}^{L}\log p(x_i\mid x_{i-1})\right]
-\Efwd\!\left[\sum_{i=1}^{L}\log q_\phi(x_{i-1}\mid x_i)\right] \\
&= \E_{x_L\sim p_L(\cdot\mid x_0)}\!\left[-\log \pprior(x_L)\right] \\
&\quad
+\E_{x_1\sim p_1(\cdot\mid x_0)}
\!\left[\log p_1(x_1\mid x_0)-\log q_\phi(x_0\mid x_1)\right] \\
&\quad
+\sum_{i=2}^{L}
\E_{(x_{i-1},x_i)\sim p(\cdot\mid x_0)}
\!\left[\log p(x_i\mid x_{i-1})-\log q_\phi(x_{i-1}\mid x_i)\right].
\end{align*}
The middle bracket equals
$\E_{x_1\sim p_1(\cdot\mid x_0)}[\log p_1(x_1\mid x_0)]+\Lrecon(x_0;\phi)$;
the first part will be cancelled by the telescoping sum below.
For each $i\ge 2$, apply Bayes' rule on the forward chain conditioned on~$x_0$:
\[
\log p(x_i\mid x_{i-1})
=\log p(x_{i-1}\mid x_i,x_0)+\log p_i(x_i\mid x_0)-\log p_{i-1}(x_{i-1}\mid x_0).
\]
Hence
\begin{align*}
&\E_{(x_{i-1},x_i)\sim p(\cdot\mid x_0)}
\!\left[\log p(x_i\mid x_{i-1})-\log q_\phi(x_{i-1}\mid x_i)\right] \\
&\quad=
\E_{x_i\sim p_i(\cdot\mid x_0)}
\!\left[\KL{p(x_{i-1}\mid x_i,x_0)}{q_\phi(x_{i-1}\mid x_i)}\right] \\
&\quad\quad
+\E_{x_i\sim p_i(\cdot\mid x_0)}\!\left[\log p_i(x_i\mid x_0)\right]
-\E_{x_{i-1}\sim p_{i-1}(\cdot\mid x_0)}
\!\left[\log p_{i-1}(x_{i-1}\mid x_0)\right].
\end{align*}
Summing the last two terms from $i=2$ to $L$ telescopes:
\begin{multline*}
\sum_{i=2}^{L}\Big(
\E_{x_i\sim p_i(\cdot\mid x_0)}\!\left[\log p_i(x_i\mid x_0)\right]
-\E_{x_{i-1}\sim p_{i-1}(\cdot\mid x_0)}
\!\left[\log p_{i-1}(x_{i-1}\mid x_0)\right]
\Big) \\
=\E_{x_L\sim p_L(\cdot\mid x_0)}\!\left[\log p_L(x_L\mid x_0)\right]
-\E_{x_1\sim p_1(\cdot\mid x_0)}\!\left[\log p_1(x_1\mid x_0)\right].
\end{multline*}
Adding the $i=1$ bracket (which contained the extra
$\E_{x_1\sim p_1(\cdot\mid x_0)}[\log p_1(x_1\mid x_0)]$ term) to the lower
end of the telescope cancels that term, leaving
\begin{align*}
\Lbound(x_0;\phi)
&= \E_{x_L\sim p_L(\cdot\mid x_0)}
\!\left[\log p_L(x_L\mid x_0)-\log \pprior(x_L)\right] \\
&\quad
+\Lrecon(x_0;\phi)
+\sum_{i=2}^{L}
\E_{x_i\sim p_i(\cdot\mid x_0)}
\!\left[\KL{p(x_{i-1}\mid x_i,x_0)}{q_\phi(x_{i-1}\mid x_i)}\right].
\end{align*}
The first bracket is $\Lprior(x_0)$ and the sum is $\Ldiff(x_0;\phi)$, hence
\[
-\log q_\phi(x_0)\le \Lprior(x_0)+\Lrecon(x_0;\phi)+\Ldiff(x_0;\phi).
\]
Taking $\E_{x_0\sim\pdata}[\cdot]$ and using the Gaussian KL identity
$\KL{p(x_{i-1}\mid x_i,x_0)}{q_\phi(x_{i-1}\mid x_i)}
=\frac{1}{2\sigma^2(i)}\|\mu(x_i,x_0,i)-\mu_\phi(x_i,i)\|_2^2$
recovers the mean-matching objective $\Ldiff(\phi)$ of Section~\ref{sec:reverse}.
\end{proof}

\begin{remark}{Role of Three Terms}{ddpm-elbo-remark}
\begin{itemize}[leftmargin=1.2em, itemsep=2pt]
    \item $\Lprior$: controlled by the noise schedule so that $p_L(\cdot\mid x_0)\approx\pprior$ and $\bpMarg{L}\approx \cN(0,\mathbf{I})$.
    \item $\Lrecon$: likelihood at the first reverse step.
    \item $\Ldiff$: matches learned reverse transitions $q_\phi(x_{i-1}\mid x_i)$ to forward posteriors $p(x_{i-1}\mid x_i,x_0)$.
\end{itemize}
\end{remark}

% ============================================================
\section{Practical Choices of Prediction and Loss}\label{sec:practical}
% ============================================================

\subsection{\texorpdfstring{$\varepsilon$}{epsilon}-Prediction Parameterization}

From the forward closed form,
\[
x_i=\balpha_i x_0+\sqrt{1-\balpha_i^2}\,\varepsilon,\qquad
\varepsilon\sim\cN(0,\mathbf{I}),
\]
the posterior mean can be rewritten as
\[
\mu(x_i,x_0,i)
=\frac{1}{\alpha_i}\!\left(
x_i-\frac{\beta_i}{\sqrt{1-\balpha_i^2}}\,\varepsilon
\right).
\]
This motivates parameterizing
\[
\mu_\phi(x_i,i)
=\frac{1}{\alpha_i}\!\left(
x_i-\frac{\beta_i}{\sqrt{1-\balpha_i^2}}\,\varepsilon_\phi(x_i,i)
\right),
\]
where $\varepsilon_\phi$ directly predicts injected noise.

\begin{remark}{Noise-Detection View}{epsilon-view}
The network is trained as a noise detector: infer the forward noise from a corrupted sample $x_i$, then use that estimate to form reverse denoising steps.
\end{remark}

\subsection{Simplified Objective}

Under the above parameterization, the mean-matching objective is equivalent (up to constants and positive time-dependent weights) to weighted noise regression:
\[
\left\|
\mu_\phi(x_i,i)-\mu(x_i,x_0,i)
\right\|_2^2
\propto
\left\|
\varepsilon_\phi(x_i,i)-\varepsilon
\right\|_2^2.
\]
The commonly used simplified loss is
\[
\Lsimple(\phi)
\triangleq
\E_{i\sim\mathrm{Uniform}\{1,\dots,L\}}\,\E_{x_0\sim\pdata}\,\E_{\varepsilon\sim\cN(0,\mathbf{I})}
\!\left[
\left\|
\varepsilon_\phi(x_i,i)-\varepsilon
\right\|_2^2
\right],
\]
where
\[
x_i=\balpha_i x_0+\sqrt{1-\balpha_i^2}\,\varepsilon.
\]

Because this is a least-squares problem, the optimal predictor at each $(x_i,i)$ is the conditional mean:
\[
\varepsilon_\phi^{*}(x_i,i)
=\E_{x_0\sim \pdata(\cdot\mid x_i)}\!\left[\frac{x_i-\balpha_i x_0}{\sqrt{1-\balpha_i^2}}\right],
\qquad x_i\sim \bpMarg{i},
\]
equivalently $\E_{x_0\sim\pdata,\,\varepsilon\sim\cN(0,\mathbf{I})}[\varepsilon\mid x_i]$ under the forward sampling
$x_i=\balpha_i x_0+\sqrt{1-\balpha_i^2}\,\varepsilon$.
Even though a single $x_i$ does not uniquely determine $x_0$, averaging over $\pdata(x_0\mid x_i)$ yields the MSE-optimal noise estimate.

\subsection{\texorpdfstring{$x_0$}{x0}-Prediction Parameterization}

An equivalent choice is to predict a clean sample directly:
\[
x_\phi(x_i,i)\approx x_0.
\]
Using Proposition~\ref{prop:reverse-posterior}, replace $x_0$ in the posterior mean by $x_\phi(x_i,i)$:
\[
\mu_\phi(x_i,i)
=\frac{\alpha_i-\balpha_i^2/\alpha_i}{1-\balpha_i^2}\,x_i
+\frac{(1-\alpha_i^2)\balpha_i/\alpha_i}{1-\balpha_i^2}\,x_\phi(x_i,i).
\]
Therefore
\[
\mu_\phi(x_i,i)-\mu(x_i,x_0,i)
=\frac{(1-\alpha_i^2)\balpha_i/\alpha_i}{1-\balpha_i^2}
\big(x_\phi(x_i,i)-x_0\big),
\]
and the diffusion KL term is equivalent to weighted clean-sample regression:
\[
\Ldiff(\phi)
=\sum_{i=2}^{L} w_i\,
\E_{x_0\sim\pdata}\,\E_{x_i\sim p_i(\cdot\mid x_0)}
\!\left[\|x_\phi(x_i,i)-x_0\|_2^2\right],
\]
for positive weights
\[
w_i
=\frac{1}{2\sigma^2(i)}
\left(
\frac{(1-\alpha_i^2)\balpha_i/\alpha_i}{1-\balpha_i^2}
\right)^2.
\]
Hence the least-squares optimal predictor is
\[
x_\phi^{*}(x_i,i)=\E[x_0\mid x_i].
\]

\subsection{Equivalence of Prediction Targets}

The targets $(\mu,x_0,\varepsilon)$ are linearly equivalent. From
\[
x_i=\balpha_i x_0+\sqrt{1-\balpha_i^2}\,\varepsilon,
\]
we have
\[
x_0=\frac{x_i-\sqrt{1-\balpha_i^2}\,\varepsilon}{\balpha_i},
\qquad
\varepsilon=\frac{x_i-\balpha_i x_0}{\sqrt{1-\balpha_i^2}}.
\]
Hence prediction heads are bijectively related:
\[
x_\phi(x_i,i)
=\frac{x_i-\sqrt{1-\balpha_i^2}\,\varepsilon_\phi(x_i,i)}{\balpha_i},
\qquad
\varepsilon_\phi(x_i,i)
=\frac{x_i-\balpha_i x_\phi(x_i,i)}{\sqrt{1-\balpha_i^2}}.
\]
So one may train with either $x_0$-prediction or $\varepsilon$-prediction.

% ============================================================
\section{Sampling from a Trained DDPM}\label{sec:sampling}
% ============================================================

After training $q_\phi(x_{i-1}\mid x_i)$, generate samples by reverse-time recursion.
Start from
\[
x_L\sim \pprior=\cN(0,\mathbf{I}),
\]
then for $i=L,L-1,\dots,1$, sample
\[
x_{i-1}
=\mu_\phi(x_i,i)+\sigma(i)\,z_i,\qquad
z_i\sim\cN(0,\mathbf{I}),
\]
with $z_1\equiv 0$ for the last step.

\paragraph{(a) $\varepsilon$-prediction form.}
Using $\mu_\phi(x_i,i)=\frac{1}{\alpha_i}\!\left(x_i-\frac{\beta_i}{\sqrt{1-\balpha_i^2}}\,\varepsilon_\phi(x_i,i)\right)$,
\[
x_{i-1}
=\frac{1}{\alpha_i}\!\left(
x_i-\frac{\beta_i}{\sqrt{1-\balpha_i^2}}\,\varepsilon_\phi(x_i,i)
\right)+\sigma(i)\,z_i.
\]

\paragraph{(b) $x_0$-prediction form.}
Using $x_\phi(x_i,i)$ instead of $\varepsilon_\phi$,
\[
x_{i-1}
=\frac{\alpha_i-\balpha_i^2/\alpha_i}{1-\balpha_i^2}\,x_i
+\frac{(1-\alpha_i^2)\balpha_i/\alpha_i}{1-\balpha_i^2}\,x_\phi(x_i,i)
+\sigma(i)\,z_i.
\]
Both forms are equivalent under the linear map in the previous subsection.

\begin{remark}{Denoise-Then-Re-noise View}{denoise-renoise-view}
Each reverse step can be interpreted as:
\begin{itemize}[leftmargin=1.2em, itemsep=2pt]
    \item \textbf{Denoise:} use the network prediction from $x_i$ (either $x_\phi$ or $\varepsilon_\phi$) to estimate the clean structure;
    \item \textbf{Re-noise:} inject calibrated Gaussian uncertainty through $\sigma(i)z_i$ to obtain $x_{i-1}$.
\end{itemize}
This is exactly the DDPM posterior-sampling mechanism under the learned mean.
\end{remark}

\end{document}
